\section{Appendix: Scope and Qualifications}
This is a study of the universal 1962 Roman Missal formulary for Our Lord Jesus Christ the King on 25 October 2026. It covers all eleven active proper elements and retains the two inactive seasonal votive witnesses in the preceding appendix. Unspecified local calendars and privileged local occurrences are outside its scope. It is a study aid, not an official liturgical book or a substitute for the governing rubrics.

The full Latin follows the independently collated 1962 controller, printed pp. 697--698 and 291. Its preexisting proper spans were compared with \emph{Acta Apostolicae Sedis} XVII (1925), pp. 666--668; the common Latin framing has an independently checked 1862 Pustet antecedent. The English reproduces the continuous Lasance--Walsh witness, revised edition of 1945, in the registered later printing, pp. 1257--1260 and 813--814, with common framing from pp. 757, 772 and 819. It is historical lay-book English, not an approved English translation for celebration. The historical proper Latin and specified English body pages have separately recorded United States public-domain bases; that finding does not clear entire facsimiles or excluded later inserts. Third-party wording is not offered under the project's license.

The source record preserves the English witness's printed \emph{familes}, repeated \emph{the} in the Epistle heading, \emph{The the Gentiles} in the Offertory, and punctuation peculiarities. Lasance's Latin reporting verb at Pilate's second question differs from the 1962 controller; the latter governs the Latin printed here. Its older Sunday-commemoration instruction is excluded. The Introit's psalm petition is numbered 71:1 by the Missal but falls at 71:2 in the inspected American-1899 Douay delivery, whose title occupies verse 1. The Communion selects clauses from both Psalm 28:10 and 11. The Secret's final dialogue and Preface's opening represent one shared ceremonial transition, displayed at both textual boundaries, not repeated in celebration.

The full inspection outcomes, textual expansions and evidence ownership are in \texttt{propers/verified.md} and \texttt{research/scope.md}.

The two whole-formulary interpretations are editorial syntheses grounded in the checked arguments of the named Fathers and saints. Their sources expound particular biblical texts, not this feast Mass as a whole. The reception inquiry checked the major accessible witnesses recorded in the research matrix; it did not establish exhaustive coverage. No direct later saintly exposition of Apocalypse 5:12 was verified. The received Victorinus text has attribution and interpolation qualifications and is contextual support only. The NPNF rendition of Augustine's Psalm 71 contains printed omissions; it is not represented as a complete examination of the Latin exposition. Modern editorial notes in NPNF remain distinct from the Fathers' arguments. Pius XI supplies the institution and doctrinal context of the feast, not an eligible saint counted among the interpretations' carriers.

The historical appendix preserves the canonical chronology's alternatives and unresolved states, alongside the reviewed modern Catholic scholarly qualifications for John and Colossians. Chronology dates its stated subjects; it does not date creation, divine kingship or the eternal kingdom. The study carries no claim of ecclesiastical approval or human specialist review.


\section{References}
\begin{itemize}[leftmargin=1.2em,itemsep=.55em]
\item \textbf{Missal and historical text witnesses.}
\emph{Missale Romanum}, Vatican typical edition, 1962, pp. 697--698, nos. 3938--3949; p. 291, no. 1080; \emph{Rubricae generales}, 17(d), 111(a), 115(c). \href{https://media.churchmusicassociation.org/pdf/missale62.pdf}{CMAA facsimile}. \emph{Acta Apostolicae Sedis} XVII (1925), fasc. 15, pp. 666--668. Francis Xavier Lasance and Francis Augustine Walsh, \emph{The New Roman Missal}, revised edition 1945, registered later printing, body pp. 1257--1260, Preface pp. 813--814, common framing pp. 757, 772, 819. \emph{Missale Romanum}, Pustet, Ratisbon, 1862, common formulas identified in the source record. These historical witnesses retain their own rights status.
\item \textbf{Biblical context.}
Douay--Rheims--Challoner, American edition 1899, registered verse deliveries. Used for complete literary contexts and comparison with liturgical adaptations; the displayed proper English is Lasance's continuous wording. Exact biblical loci appear with each discussion.
\item \textbf{Augustine of Hippo.}
\emph{Tractates on the Gospel of John}, 115.1--5, in \emph{Nicene and Post-Nicene Fathers}, first series, VII, historical English, CCEL delivery. \emph{Expositions of the Psalms}, Vulgate Psalms 2, 28 and 71, in NPNF, first series, VIII; the latter two appear under English headings 29 and 72. Exact sections used are cited in the prose. \href{https://www.ccel.org/ccel/schaff/npnf107.html}{NPNF I.7}; \href{https://www.ccel.org/ccel/schaff/npnf108.html}{NPNF I.8}.
\item \textbf{Jerome.}
\emph{Commentaria in Danielem}, VII, on 7:13--14,17--18,27, PL 25, cols. 670--672, retained historical Latin transcription. \href{https://la.wikisource.org/wiki/Commentaria_in_Danielem_(Hieronymus)}{Latin witness}. Generated transcription paragraph markers are not native loci.
\item \textbf{John Chrysostom.}
\emph{Homilies on Colossians}, II on 1:12--15 and III, in NPNF, first series, XIII, historical English, \href{https://www.ccel.org/ccel/schaff/npnf113.html}{CCEL delivery}. Textual variants and nineteenth-century editorial commentary do not replace the Missal's appointed wording.
\item \textbf{Robert Bellarmine.}
\emph{A Commentary on the Book of Psalms}, John O'Sullivan translation, Psalms 2, 28 and 71, retained historical-text derivative of the \href{https://www.ecatholic2000.com/bellarmine/commentary-on-psalms.shtml}{eCatholic2000 delivery}. Numbered comment 10 under Psalm 71 treats biblical verse 11.
\item \textbf{Received Victorinus commentary.}
\emph{Commentary on the Apocalypse}, ch. 1 on vv. 4--7 and ch. 5, ANF VII, historical English with the edition's attribution and interpolation cautions. \href{https://www.ccel.org/ccel/schaff/anf07.html}{ANF VII}.
\item \textbf{Feast context.}
Pius XI, \emph{Quas primas}, 11 December 1925, nos. 12--20, \href{https://www.vatican.va/content/pius-xi/en/encyclicals/documents/hf_p-xi_enc_11121925_quas-primas.html}{official English delivery}; used through original analysis, not reproduced translation. FSSP Province de France, \emph{Ordo du mois}, October 2026 state, 25 October entry, as identified in the source record.
\item \textbf{Historical orientation.}
\emph{The Catholic Encyclopedia}: Corbett, \emph{King David}, IV, opening; Gigot, \emph{Book of Daniel}, IV, authorship; Howlett, \emph{Biblical Chronology}, III, Passion discussion; \emph{Gospel of St John}, VIII, circumstances of composition; \emph{Epistle to the Colossians}, IV, captivity; \emph{St Paul}, XI, chronology list; \emph{Apocalypse}, I, time and place; \emph{New Testament}, XIV, dating. Historical article-text witnesses and exact locators are bound in the research records.
\item \textbf{Modern Catholic historical introductions.}
USCCB, NABRE introductions to \href{https://bible.usccb.org/bible/psalms/0}{Psalms}, superscriptions and dating; \href{https://bible.usccb.org/bible/john/0}{John}, paragraphs 5--7 and 10; \href{https://bible.usccb.org/bible/colossians/0?bk=Colossians&ch=}{Colossians}, paragraphs 2 and 5. Reviewed October 2026 deliveries; original historical analysis only, with no reproduction of protected page bodies.
\end{itemize}
